\documentclass[12pt,a4paper]{article}
\usepackage[utf8]{inputenc}
\usepackage[brazil]{babel}
\usepackage[margin=2cm]{geometry}
\usepackage{booktabs}
\usepackage{tabularx}
\usepackage{xcolor}
\usepackage{enumitem}
\usepackage{framed}
\usepackage{fancyhdr}
\usepackage{array}
\usepackage{amssymb}
\usepackage{graphicx}
\usepackage{multirow}
\usepackage{tikz}

\sloppy
\definecolor{ufpiblue}{RGB}{0,51,102}
\definecolor{farmaigreen}{RGB}{14,165,122}
\definecolor{evalamber}{RGB}{234,179,8}
\definecolor{evalred}{RGB}{239,68,68}
\definecolor{accentpro}{RGB}{96,165,250}

\pagestyle{fancy}
\fancyhf{}
\fancyhead[L]{\small\textcolor{ufpiblue}{\textbf{UFPI}} -- Dept. Computação}
\fancyhead[R]{\small Questionário FarmaIA}
\fancyfoot[C]{\small Página \thepage}
\renewcommand{\headrulewidth}{0.4pt}

% Reusable command for concordance row
\newcommand{\concrow}[1]{%
  #1 & $\square$ & $\square$ & $\square$ & \\[10pt]%
}

\begin{document}

% ── CAPA ──
\begin{titlepage}
\centering
\vspace*{1cm}
{\Large\textcolor{ufpiblue}{\textbf{Universidade Federal do Piauí}}}\\[2mm]
{\large Departamento de Computação}\\[1cm]
\rule{\textwidth}{1pt}\\[5mm]
{\LARGE\textbf{Questionário de Avaliação}}\\[3mm]
{\Large\textcolor{farmaigreen}{\textbf{FarmaIA}}: Validação dos Juízes Automáticos}\\[5mm]
\rule{\textwidth}{1pt}
\vspace{1.2cm}

\begin{flushleft}
\large
\textbf{Objetivo:} Verificar se os juízes automáticos do FarmaIA
avaliam as respostas de forma compatível com o julgamento de
estudantes de Farmácia.

\vspace{6mm}
\textbf{Como funciona:}
\begin{enumerate}[leftmargin=*,itemsep=4pt]
  \item A interface apresenta uma sessão de chat com o FarmaIA.
    No painel esquerdo você verá cada par pergunta/resposta, junto
    com os detalhes da \textbf{avaliação automática MCP} (nota
    agregada, juízes gerais e juízes de tópico).
  \item Para cada resposta do sistema, clique em
    \textbf{``Iniciar Avaliação''} para abrir o formulário de
    concordância no painel direito.
  \item Para cada juiz (geral e de tópico), indique se você
    \textbf{concorda}, \textbf{concorda parcialmente} ou
    \textbf{discorda} da avaliação automática.
  \item Se concordar parcialmente ou discordar, escreva
    brevemente o motivo.
  \item Você também pode enviar novas perguntas no chat para
    gerar respostas adicionais com avaliação automática.
  \item Ao final, preencha a \textbf{Avaliação Geral} e clique
    em \textbf{``Enviar Avaliação Completa''}.
\end{enumerate}

\vspace{6mm}
\textbf{Tempo estimado:} 40--60 minutos. \hfill
\textbf{Questionário anônimo.}
\end{flushleft}

\vfill
\rule{\textwidth}{0.5pt}
\vspace{3mm}

\noindent
\textbf{Código do avaliador:} \textit{gerado automaticamente pela interface} \hfill
\textbf{Data:} \framebox[4cm]{\rule{0pt}{7mm}}

\vspace{4mm}
\noindent\textbf{Semestre atual:} \framebox[4cm]{\rule{0pt}{7mm}} \hfill
\textbf{Cursou Farmacologia?} \hspace{3mm}
$\square$~Sim \hspace{5mm} $\square$~Não
\end{titlepage}

% ════════════════════════════════════════════════════════════
% ESTRUTURA DA INTERFACE
% ════════════════════════════════════════════════════════════
\section*{Estrutura da Interface de Avaliação}

A interface é dividida em dois painéis:

\begin{itemize}[leftmargin=*,itemsep=4pt]
  \item \textbf{Painel Esquerdo --- Conversação \& Juízes:}
    Exibe todas as mensagens do chat (perguntas do usuário e
    respostas do FarmaIA). Cada resposta possui um painel expansível
    ``Avaliação MCP'' com a nota agregada e os detalhes de cada juiz.
  \item \textbf{Painel Direito --- Sua Avaliação:}
    Formulário de concordância que aparece ao clicar em
    ``Iniciar Avaliação'' em uma resposta.
\end{itemize}

\subsection*{Juízes Gerais (4)}

Para cada resposta, o sistema executa os seguintes juízes gerais,
cada um atribuindo uma nota de 0 a 100:

\begin{center}
\small
\begin{tabularx}{\textwidth}{@{}l l X@{}}
\toprule
\textbf{Juiz} & \textbf{Identificador} & \textbf{O que avalia} \\
\midrule
🛡️ Segurança & \texttt{safety\_judge} &
  Se a resposta é segura, com dosagens corretas e alertas
  necessários. \\[6pt]
⭐ Qualidade & \texttt{quality\_judge} &
  Se a resposta é completa, precisa e relevante para a
  pergunta. \\[6pt]
📎 Fundamentação & \texttt{source\_judge} &
  Se as afirmações são rastreáveis à bula (inclui lista de
  \textit{claims} classificados como \textit{Supported},
  \textit{Unsupported} ou \textit{Partially Supported}). \\[6pt]
📐 Formato & \texttt{format\_judge} &
  Se a organização, linguagem e extensão são adequadas. \\
\bottomrule
\end{tabularx}
\end{center}

Cada juiz geral fornece justificativas estruturadas por critério
(ex.: \textit{Precisão da dosagem}, \textit{Alertas de segurança},
etc.) ao invés de uma justificativa única em texto livre.

\subsection*{Juízes de Tópico (até 3)}

Quando o sistema detecta tópicos específicos na resposta, juízes
adicionais são ativados com uma nota de cobertura
(\texttt{coverage\_score}):

\begin{center}
\small
\begin{tabularx}{\textwidth}{@{}l l X@{}}
\toprule
\textbf{Juiz de Tópico} & \textbf{Identificador} & \textbf{O que avalia} \\
\midrule
💊 Posologia & \texttt{posologia} &
  Cobertura de informações de dosagem, frequência e via de
  administração. \\[6pt]
⚠️ Contraindicações & \texttt{contraindicacoes} &
  Cobertura de contraindicações e interações
  medicamentosas. \\[6pt]
🔬 Reações Adversas & \texttt{reacoes\_adversas} &
  Cobertura de efeitos colaterais e reações adversas. \\
\bottomrule
\end{tabularx}
\end{center}

Cada juiz de tópico lista as \textbf{perguntas respondidas}
(✅) e as \textbf{perguntas não respondidas} (❌) pela resposta.

\newpage

% ════════════════════════════════════════════════════════════
% MODELO DE AVALIAÇÃO POR RESPOSTA
% ════════════════════════════════════════════════════════════
\section*{Modelo de Avaliação por Resposta \hfill (repetido para cada resposta)}

Para cada resposta do FarmaIA com avaliação automática, o avaliador
vê os seguintes dados no painel esquerdo e preenche o formulário
de concordância no painel direito.

\vspace{3mm}
\subsection*{Painel Esquerdo: Dados da Resposta}

\begin{framed}
\noindent\textbf{Pergunta do Usuário:}\\[2pt]
\textit{[Texto da pergunta original]}
\end{framed}

\begin{framed}
\noindent\textbf{Resposta do FarmaIA:}\\[2pt]
\textit{[Texto da resposta gerada pelo sistema, renderizado em Markdown]}
\end{framed}

\begin{framed}
\noindent\textbf{Avaliação MCP} --- Nota agregada: \framebox[2cm]{\rule{0pt}{5mm}}/100\\[4pt]

\noindent\textit{Juízes gerais:}\\[2pt]
Para cada juiz geral, exibe-se:
\begin{itemize}[leftmargin=12pt,itemsep=2pt]
  \item Nome do juiz e nota (badge colorido: {\color{farmaigreen}$\geq$80},
    {\color{evalamber}60--79}, {\color{evalred}$<$60})
  \item Justificativas estruturadas (uma por critério avaliado)
  \item Para o juiz de Fundamentação: lista de \textit{claims}
    extraídos da resposta, cada um classificado como
    {\color{farmaigreen}\textit{Supported}},
    {\color{evalamber}\textit{Partially Supported}} ou
    {\color{evalred}\textit{Unsupported}}
\end{itemize}

\noindent\textit{Juízes de tópico (quando detectados):}\\[2pt]
Para cada tópico detectado:
\begin{itemize}[leftmargin=12pt,itemsep=2pt]
  \item Nome do tópico e nota de cobertura (badge colorido)
  \item Lista de perguntas respondidas (✅) e não respondidas (❌)
\end{itemize}
\end{framed}

\vspace{3mm}
\subsection*{Painel Direito: Formulário de Concordância}

Para \textbf{cada juiz} (gerais e de tópico), o avaliador responde:

\vspace{2mm}
\small
\begin{tabularx}{\textwidth}{@{}X c c c p{3cm}@{}}
\toprule
\textbf{Juiz (com nota)} &
\rotatebox{90}{\textbf{Concordo}} &
\rotatebox{90}{\textbf{Parcial}} &
\rotatebox{90}{\textbf{Discordo}} &
\textbf{Se parcial/discorda, por quê?} \\
\midrule
\concrow{🛡️ Segurança \hfill \framebox[1.5cm]{\rule{0pt}{4mm}}/100}
\concrow{⭐ Qualidade \hfill \framebox[1.5cm]{\rule{0pt}{4mm}}/100}
\concrow{📎 Fundamentação \hfill \framebox[1.5cm]{\rule{0pt}{4mm}}/100}
\concrow{📐 Formato \hfill \framebox[1.5cm]{\rule{0pt}{4mm}}/100}
\midrule
\multicolumn{5}{@{}l}{\textit{Juízes de tópico (se detectados):}} \\[4pt]
\concrow{💊 Posologia (Tópico) \hfill \framebox[1.5cm]{\rule{0pt}{4mm}}/100}
\concrow{⚠️ Contraindicações (Tópico) \hfill \framebox[1.5cm]{\rule{0pt}{4mm}}/100}
\concrow{🔬 Reações Adversas (Tópico) \hfill \framebox[1.5cm]{\rule{0pt}{4mm}}/100}
\bottomrule
\end{tabularx}
\normalsize

\vspace{3mm}
\noindent\textit{Nota: O campo ``Por quê?'' aparece automaticamente
quando o avaliador seleciona ``Parcial'' ou ``Discordo''.
Na versão impressa, preencher apenas quando aplicável.}

\newpage

% ════════════════════════════════════════════════════════════
% FUNCIONALIDADE DE CHAT LIVRE
% ════════════════════════════════════════════════════════════
\section*{Chat Livre durante a Avaliação}

O avaliador pode enviar novas perguntas ao FarmaIA diretamente
pela interface de avaliação, usando o campo de texto na parte
inferior do painel esquerdo. Cada nova resposta recebe
automaticamente uma avaliação dos juízes e pode ser avaliada
pelo concordância da mesma forma.

\vspace{4mm}
\begin{framed}
\noindent\textbf{Campo de chat:}
\textit{``Faça uma nova pergunta\ldots''}\\[4pt]
\noindent O sistema processa a pergunta, gera a resposta via
FarmaIA, e executa os juízes automaticamente. O avaliador pode
então clicar em ``Iniciar Avaliação'' na nova resposta.
\end{framed}

\newpage

% ════════════════════════════════════════════════════════════
% AVALIAÇÃO GERAL
% ════════════════════════════════════════════════════════════
\section*{Avaliação Geral do Sistema}
\rule{\textwidth}{0.5pt}
\vspace{4mm}

\noindent\textbf{1. Os juízes avaliaram de forma justa?}

\vspace{2mm}
\noindent
$\square$~Sim \hfill
$\square$~Parcialmente \hfill
$\square$~Não

\vspace{5mm}
\noindent\textbf{2. Qual juiz mais distante do seu julgamento?}

\vspace{2mm}
\noindent
$\square$~🛡️ Segurança \hfill
$\square$~⭐ Qualidade \hfill
$\square$~📎 Fundamentação \hfill
$\square$~📐 Formato

\vspace{5mm}
\noindent\textbf{3. Tendência de notas dos juízes?}

\vspace{2mm}
\noindent
$\square$~Muito altas (lenientes) \hfill
$\square$~Adequadas \hfill
$\square$~Muito baixas (rigorosos)

\vspace{5mm}
\noindent\textbf{4. Comentários:}
\vspace{2cm}
\noindent\rule{\textwidth}{0.3pt}
\vspace{5mm}
\noindent\rule{\textwidth}{0.3pt}
\vspace{5mm}
\noindent\rule{\textwidth}{0.3pt}

\vspace{4mm}
\begin{center}
\fbox{\parbox{0.85\textwidth}{\centering\vspace{2mm}
\textbf{Enviar Avaliação Completa}\\[2pt]
\small Na interface web, o botão submete todos os dados de
concordância e a avaliação geral ao servidor.
\vspace{2mm}}}
\end{center}

\vfill
\begin{center}
\rule{0.5\textwidth}{0.5pt}\\[4mm]
\textbf{Obrigado pela participação!}\\[2mm]
\small TCC de Paulo Eduardo B. do Vale -- Orientador: Prof.~Pedro
Santos Neto -- UFPI
\end{center}

\end{document}
